% !TeX program = lualatex
% Compile twice: lualatex 109.tex
% All references and prose are embedded; no BibTeX database or external inputs.
\documentclass[11pt,a4paper]{article}
\usepackage[margin=24mm,headheight=14pt]{geometry}
\usepackage{fontspec}
\setmainfont{Latin Modern Roman}
\setsansfont{Latin Modern Sans}
\setmonofont{Latin Modern Mono}
\usepackage{amsmath,amssymb,mathtools}
\usepackage{unicode-math}
\setmathfont{Latin Modern Math}
\usepackage{microtype}
\usepackage{enumitem,xurl,needspace}
\usepackage[dvipsnames]{xcolor}
\usepackage{hyperref}
\hypersetup{unicode=true,colorlinks=true,linkcolor=MidnightBlue,urlcolor=MidnightBlue,
 pdftitle={Research attempt: Problem 109},pdfauthor={Research notebook},bookmarksnumbered=true}
\usepackage{bookmark}
\usepackage{tocloft}
\setlength{\cftsecnumwidth}{3em}
\cftsetpnumwidth{2.5em}
\cftsetrmarg{4em}
\usepackage{titlesec}
\titleformat{\section}{\Large\bfseries\raggedright}{\thesection}{0.7em}{}
\usepackage{fancyhdr}
\pagestyle{fancy}\fancyhf{}
\fancyhead[L]{\small\sffamily Open Problems Math | Research notebook}
\fancyhead[R]{\small Problem 109 | 27 September 2026}
\fancyfoot[C]{\thepage}
\setlength{\parindent}{0pt}
\setlength{\parskip}{3pt plus 1pt}
\setlength{\emergencystretch}{3em}
\setcounter{tocdepth}{1}
\setcounter{secnumdepth}{1}
\setlist[itemize]{leftmargin=3em,itemsep=5pt,topsep=4pt,parsep=0pt}
\allowdisplaybreaks
\newcommand{\N}{\mathbb{N}}
\newcommand{\Z}{\mathbb{Z}}
\newcommand{\Q}{\mathbb{Q}}
\newcommand{\R}{\mathbb{R}}
\newcommand{\C}{\mathbb{C}}
\newcommand{\F}{\mathbb{F}}
\newcommand{\A}{\mathbb{A}}
\newcommand{\PP}{\mathbb P}
\newcommand{\E}{\mathbb{E}}
\newcommand{\eps}{\varepsilon}
\newcommand{\dd}{\,\mathrm d}
\newcommand{\sref}[2]{\hyperlink{source:#1:#2}{[S#2]}}
\newcommand{\eref}[2]{\hyperlink{extra:#1:#2}{[E#2]}}
% Turn the next switch off for a shorter reading copy. The quotations preserve
% every exactTarget field, including qualifications omitted from a short summary.
\newif\ifshowcatalogue
\showcataloguetrue
\newcommand{\cataloguescope}[1]{\ifshowcatalogue
 \paragraph{Exact scope in the supplied catalogue.}
 \begin{quote}\small #1\end{quote}\fi}
\usepackage{etoolbox}
\AtBeginEnvironment{itemize}{\small}
\numberwithin{equation}{section}
% Standalone export. Original section 109 preserved verbatim outside marked additions.
\begin{document}
\setcounter{section}{108}
% ================================================================
% problemId: problem.global-regularity-of-the-3d-incompressible-magnetohydrodynamic-equations
% source releaseRank: 109; source releaseStatus: open_with_solved_subcases
% exactTarget SHA-256: e0d2e073bcfe3d57bb1d9e4c4756d374af95224d026fc7f11116570dad2c376d
\clearpage
\section{\texorpdfstring{Global Regularity of the 3D Incompressible Magnetohydrodynamic Equations}{Global Regularity of the 3D Incompressible Magnetohydrodynamic Equations}}\label{109}
\subsection{Definitions and mathematical statement}
For velocity $u$, magnetic field $B$, and total pressure $p$, the standard viscous and resistive incompressible MHD equations are
\[
 \begin{aligned}
 \partial_tu+(u\cdot\nabla)u+\nabla p&=\nu\Delta u+(B\cdot\nabla)B,\\
 \partial_tB+(u\cdot\nabla)B&=\eta\Delta B+(B\cdot\nabla)u,\\
 \nabla\cdot u=\nabla\cdot B&=0,
 \end{aligned}\qquad x\in\R^3,\quad\nu,\eta>0.
\]
The target is unique global smooth evolution for every allowed smooth divergence-free initial pair. The catalogue's phrase ``sufficiently smooth'' does not fix a Sobolev threshold or decay condition; these must match the cited formulation. Pressure is understood up to addition of a function of time. Small-data regularity does not cover arbitrary data.
\par\smallskip\textit{Formulation and concept references:} \sref{109}{1}.
\cataloguescope{Prove or disprove that every sufficiently smooth divergence-free initial velocity and magnetic field on R\textasciicircum{}3 generates a unique global smooth solution of the standard incompressible three-dimensional MHD system with positive viscosity and magnetic diffusivity.}
\subsection{Short English statement}
Do viscosity and magnetic diffusion prevent singularities for every allowed smooth three-dimensional incompressible magnetofluid flow?
% BEGIN RESEARCH ADDITION 109
\subsection{Research attempt}

\paragraph{Current frontier and source audit.}
The publisher's first-page preview of \sref{109}{1}, printed p.~14417,
explicitly uses divergence-free data in $H^s(\R^3)$, $s\geq3$, and separates
local strong well-posedness from global weak existence. Its subsequent
normalization $\nu=\eta=1$ is a simplification of that paper, not a license
to impose equal diffusivities on the present target. Here the initial data
are smooth and belong to that finite-energy Sobolev setting; all estimates
are first made on their classical interval of existence. The original
statement is left unchanged. \sref{109}{1}

The 2024 paper obtains conditional, logarithmically improved Besov-space
criteria. The same authors' May 2026 paper obtains further conditional
criteria involving selected velocity/magnetic derivatives and pressure.
Neither result supplies the requisite estimates from arbitrary initial
data. Only the introduction/preview of the former and the publisher's
abstract of the latter were accessible in this audit; no inaccessible
technical theorem is used below. \sref{109}{1}\,\eref{109}{1}

\paragraph{Attack route.}
There are three natural ways to exploit the magnetic coupling. A positive
cross-helicity might make the two fields align, the Elsasser variables might
leave only one dangerous advecting field, or the ordinary energy inequality
might control a critical continuation norm. The first route needs a
mechanism forcing alignment, which energy conservation does not provide.
The third route has a scaling obstruction. We pursue the second route,
keeping the unequal-diffusivity term rather than discarding it. The result
will be a one-Elsasser-field continuation estimate with a fully displayed
nonlinear cancellation, followed by tests of the missing global estimate.

\paragraph{Attempt: a one-sided enstrophy estimate.}
Put
\[
 z^+=u+B,\qquad z^-=u-B,\qquad
 a=\frac{\nu+\eta}{2},\quad b=\frac{\nu-\eta}{2},\quad
 m=\min(\nu,\eta)>0.
\]
Adding and subtracting the original equations gives, without changing the
pressure convention,
\begin{equation}\label{109:elsasser}
 \partial_tz^\pm+(z^\mp\cdot\nabla)z^\pm+\nabla p
 =a\Delta z^\pm+b\Delta z^\mp,
 \qquad \nabla\cdot z^\pm=0.
\end{equation}
Define
\[
 Y=\|\nabla z^+\|_2^2+\|\nabla z^-\|_2^2,\qquad
 D=\|\Delta z^+\|_2^2+\|\Delta z^-\|_2^2.
\]
The diffusion quadratic form obeys
\begin{equation}\label{109:diffusion}
 aD+2b\langle\Delta z^+,\Delta z^-\rangle\geq mD,
\end{equation}
since the $2\times2$ coefficient matrix has eigenvalues $\nu,\eta$.
Testing \eqref{109:elsasser} with $-\Delta z^\pm$ kills the pressure.
The two nonlinear terms, with their actual energy-identity signs, are
\[
 I_+=\int (z^-\cdot\nabla)z^+\cdot\Delta z^+,
 \qquad I_-=\int (z^+\cdot\nabla)z^-\cdot\Delta z^-.
\]
The apparent difficulty is that $I_-$ has the large field $z^+$ as
advector. Integrating first in $x_k$, and then in $x_j$, gives
\begin{align}
 I_-&=-\sum_{i,j,k}\int
       (\partial_k z^+_j)(\partial_j z^-_i)(\partial_k z^-_i)
       \notag\\
 &=\sum_{i,j,k}\int z^-_i(\partial_k z^+_j)
                         (\partial_j\partial_k z^-_i).
 \label{109:cancellation}
\end{align}
The additional term in the second integration is zero because
$\sum_j\partial_j\partial_k z^+_j=0$. Thus both nonlinear terms contain
an \emph{undifferentiated} $z^-$. This is an identity, not an assumed
alignment or a sign assertion.

For $3<q<\infty$, set $r=2q/(q-2)$ and $\theta=3/q$.
Holder and Sobolev interpolation give
\begin{align}
 |I_+|+|I_-|
 &\leq C\|z^-\|_q\|\nabla z^+\|_r D^{1/2} \notag\\
 &\leq C\|z^-\|_q Y^{(1-\theta)/2}D^{(1+\theta)/2} \notag\\
 &\leq \tfrac m2D+
 C_qm^{-(q+3)/(q-3)}\|z^-\|_q^{2q/(q-3)}Y.
 \label{109:young}
\end{align}
Here $\|\nabla^2v\|_2=\|\Delta v\|_2$ follows by Plancherel; smooth
cutoffs or Sobolev density justify the integrations on $\R^3$.
Consequently, for $p_q=2q/(q-3)$,
\begin{equation}\label{109:gronwall}
 Y'(t)+mD(t)\leq
 C_qm^{-(q+3)/(q-3)}\|z^-(t)\|_q^{p_q}Y(t).
\end{equation}
In particular, if a finite classical lifespan endpoint $T$ satisfied
\begin{equation}\label{109:missing}
 \int_0^T\|u(t)-B(t)\|_q^{p_q}\,dt<\infty
 \quad\text{for some }q>3,
\end{equation}
then $Y$ would stay bounded and $\int_0^TD<\infty$.
Interchanging $+$ and $-$ gives the corresponding assertion for $u+B$.
This is an a priori criterion, not a claim that either norm is finite.

For completeness, the continuation implication can be closed without
assuming an $H^s$ bound as an extra hypothesis. The $L^2$ energy bound
and the preceding $H^1$ bound give bounded $L^\infty_tH^1_x$ and
$L^2_tH^2_x$ norms. Writing each transport term as a divergence, the tame product bound
\[
 \|vw\|_{H^s}\leq C_s(\|v\|_\infty\|w\|_{H^s}
                          +\|w\|_\infty\|v\|_{H^s})
\]
and one integration by parts bound the higher-order nonlinear energy by
$C_s\|z\|_\infty\|z\|_{H^s}\|\nabla z\|_{H^s}$, where $\|z\|$
denotes the sum of the two fields' norms. Since
$H^2(\R^3)\hookrightarrow L^\infty(\R^3)$, we have
$\int_0^T\|z(t)\|_\infty^2\,dt<\infty$.
Young's inequality absorbs the last factor into the positive diffusion,
leaving an integrable coefficient times $\|z\|_{H^s}^2$.
Gronwall bounds each fixed $H^s$ norm up to $T$, and the local strong
theory restarts. The standard product/commutator technology is identified
in \eref{109}{2}; the one-sided cancellation above is derived here.

At the endpoint $q=3$, \eqref{109:young} instead gives
\[
 |I_+|+|I_-|\leq C\|z^-\|_3D.
\]
Thus the additional hypothesis
$\sup_{t<T}\|z^-(t)\|_3<m/(2C)$ permits direct absorption. This is a
smallness condition on an evolving critical norm; it is not an
arbitrary-data theorem.

\paragraph{Stress test: three ways the global step could fail.}
First, the unconditional energy identity is only
\[
 \frac12\frac d{dt}(\|u\|_2^2+\|B\|_2^2)
 +\nu\|\nabla u\|_2^2+\eta\|\nabla B\|_2^2=0.
\]
It provides the interpolation line $2/p+3/q=3/2$, whereas
\eqref{109:missing} lies on $2/p+3/q=1$.
The discrepancy cannot be repaired by an interpolation inequality with
these two energy norms alone. To see this, choose a nonzero smooth compactly
supported divergence-free $\phi$ and a nonzero smooth time cutoff $\chi$,
and define, for $\lambda\geq1$,
\[
 v_\lambda(t,x)=\lambda^{3/2}\chi(\lambda^2t)\phi(\lambda x).
\]
Both $\|v_\lambda\|_{L^\infty_tL^2_x}$ and
$\|\nabla v_\lambda\|_{L^2_{t,x}}$ are uniformly bounded, but
\[
 \|v_\lambda\|_{L^{p_q}_tL^q_x}
 =c_{\phi,\chi,q}\lambda^{1/2}.
\]
These are test fields, \emph{not} MHD solutions or singularity candidates.
They falsify only the proposed energy-to-critical-norm interpolation step.

Second, alignment is an invariant regime only when diffusivities permit it.
If $\nu=\eta$ and $z^-_0=0$, then
$z^-=0$, $z^+(t)=e^{\nu t\Delta}z^+_0$, and spatially constant pressure
solve the system globally. But for $\nu\ne\eta$ and $z^-_0=0$,
\eqref{109:elsasser} gives
\[
 \partial_tz^-\big|_{t=0}=b\Delta z^+_0
\]
after projecting onto divergence-free fields. A nonzero finite-energy
$z^+_0$ generally generates the supposedly absent field immediately.
Large-data equal-diffusivity alignment cannot silently be extended to the
whole parameter range.

Third, with $B_0=0$, uniqueness keeps $B=0$ and the velocity solves exactly
the three-dimensional incompressible Navier--Stokes equations. Accordingly,
a bound proving \eqref{109:missing} for \emph{every} MHD solution would also
supply a critical regularity bound for arbitrary Navier--Stokes data.
This dependency is algebraic and unavoidable, not a reason to abandon the
attempt. It identifies how much genuinely new PDE information is missing.

\paragraph{Outcome.}
\textbf{Outcome: Conditional reduction.}

For all positive $\nu,\eta$, the calculation reduces continuation to the
finite critical space-time norm of just one Elsasser field. It proves the
cancellation \eqref{109:cancellation}, keeps cross-diffusion coercive,
and supplies a scaling countertest to the most immediate attempt to deduce
the required bound from energy. No novelty is asserted for the resulting
Serrin-type criterion.

A complete positive resolution still needs an arbitrary-data mechanism
forcing \eqref{109:missing}, or an essentially different estimate strong
enough to continue the solution without it. No such mechanism, and no
finite-time singularity for the prescribed smooth data, is established here.
% END RESEARCH ADDITION 109
\subsection{Sources}
\begin{itemize}\raggedright
\item[\textnormal{[S1]}]\hypertarget{source:109:1}{} Zhang, Li, and Zhang, Logarithmically Improved Regularity Criteria for the 3D Magnetohydrodynamic Equations, Mathematical Methods in the Applied Sciences 47 (2024), introduction, p. 14417.
\url{https://www.researchgate.net/publication/381615027_Logarithmically_improved_regularity_criteria_for_the_3D_magnetohydrodynamic_equations}.
{\small\textit{Catalogue formulation source.}}
% BEGIN NEW REFERENCES 109
\item[\textnormal{[E1]}]\hypertarget{extra:109:1}{}
Honglin Zhang, Li Li, and Qingning Zhang, ``Some new regularity criteria
of the incompressible 3D MHD equations,'' \emph{Acta Mathematica Scientia}
\textbf{46} (2026), 1340 ff., DOI 10.1007/s10473-026-0315-y.
Publisher abstract consulted 27 September 2026; used only to identify the
conditional nature and scope of the recent criteria.
\url{https://link.springer.com/article/10.1007/s10473-026-0315-y}.
\item[\textnormal{[E2]}]\hypertarget{extra:109:2}{}
Tosio Kato and Gustavo Ponce, ``Commutator estimates and the Euler and
Navier--Stokes equations,'' \emph{Communications on Pure and Applied
Mathematics} \textbf{41} (1988), 891--907. Sobolev commutator estimates;
also identified in the accessible bibliography of [S1].
\url{https://doi.org/10.1002/cpa.3160410704}.
% END NEW REFERENCES 109
\end{itemize}

\end{document}
